\section{Introduction}
\label{sec:introduction}

C++ is a highly versatile programming language that provides profound support for Object-Oriented Programming (OOP) paradigms, including encapsulation, inheritance, and polymorphism. As software systems expand, developers frequently construct complex object structures, such as intricate data trees or user interface component graphs. Managing operations across these heterogeneous collections often introduces significant architectural challenges.

To address recurring design problems within OOP architectures, \textit{Design Patterns} serve as proven, standardized templates representing best practices. Among the three primary classifications of design patterns, \textbf{Behavioral Patterns} focus specifically on the effective assignment of responsibilities between objects and the intricate communication mechanisms that govern their interactions, ensuring loose coupling and high cohesion \cite{gof1994}.

The \textbf{Visitor Pattern} is a prominent behavioral design pattern whose primary intent is to separate algorithms or operations from the object structures on which they operate. In traditional object-oriented design, introducing a new behavior across a complex hierarchy of elements often leads to rigid and hard-to-maintain codebases. The Visitor pattern offers an elegant architectural approach to overcome this limitation, allowing developers to extend the functionality of a system dynamically \cite{gof1994}

The objective of this report is to present a study of the Visitor design pattern within the context of C++ development. It systematically explores the fundamental problem the pattern addresses, its structural architecture, the underlying mechanics of double dispatch, modern C++ implementation variants, and its pragmatic applications in contemporary software engineering.
